\section{Depth estimates for a merging tree}

The merging construction of \cite{Styliaris2025MPUCircuits} gives numerical
recurrences for circuit depth. The results below concern those recurrences;
they do not assert that an arbitrary tensor admits the required circuit
construction or conditioning bounds.

\begin{theorem}[Product of the conditioning costs]
    \label{thm:mpu_tree_depth_product}
    \lean{MPUCircuit.depth_le_prod_conditioning}
    \leanok
    Let $T_j,q_j$ be real sequences, let $a\in\mathbb R$ and $C\geq0$, and
    suppose $q_j\geq1$ for every nonnegative integer $j$. If
    \begin{align}
        T_0     & \leq a,                  &
        T_{j+1} & \leq q_j(2T_j+C2^{j+1}),
    \end{align}
    then for every nonnegative integer $L$,
    \begin{align}
        T_L & \leq 2^L(a+CL)\prod_{j=0}^{L-1}q_j.
    \end{align}
\end{theorem}

\begin{proof}\leanok
    Induct on $L$. The initial case is $T_0\leq a$. At the induction step,
    use $\prod_{j<L}q_j\geq1$ to bound the additive overhead by
    $C2^{L+1}\prod_{j<L}q_j$. Multiplication by $q_L\geq0$ gives
    \begin{align}
        T_{L+1}
         & \leq q_L\left(2^{L+1}(a+CL)\prod_{j<L}q_j
        +C2^{L+1}\prod_{j<L}q_j\right)               \\
         & =2^{L+1}(a+C(L+1))\prod_{j<L+1}q_j.
    \end{align}
\end{proof}

\begin{theorem}[Polynomial product bound]
    \label{thm:mpu_tree_depth_polynomial}
    \lean{MPUCircuit.depth_le_polynomial_of_prod_conditioning}
    \leanok
    Under the hypotheses of Theorem~\ref{thm:mpu_tree_depth_product}, suppose
    also $a\geq0$. Let $L,c$ be nonnegative integers. If
    $\prod_{j<L}q_j\leq (2^L)^c$, then
    \begin{align}
        T_L & \leq (a+CL)(2^L)^{c+1}.
    \end{align}
\end{theorem}

\begin{proof}\leanok
    \uses{thm:mpu_tree_depth_product}
    Substitute the assumed product bound into the preceding estimate.
    Its coefficient $2^L(a+CL)$ is nonnegative.
\end{proof}

\begin{theorem}[Constant conditioning]
    \label{thm:mpu_tree_depth_constant}
    \lean{MPUCircuit.depth_le_constant_conditioning}
    \leanok
    Let $T_j$ be a real sequence, $a\in\mathbb R$, $C\geq0$, and $q\geq1$.
    Suppose $T_0\leq a$ and $T_{j+1}\leq q(2T_j+C2^{j+1})$ for every
    nonnegative integer $j$. Then
    \begin{align}
        T_L & \leq (a+CL)(2q)^L.
    \end{align}
\end{theorem}

\begin{proof}\leanok
    \uses{thm:mpu_tree_depth_product}
    Apply the product estimate to $q_j=q$ and use $2^Lq^L=(2q)^L$.
\end{proof}

\begin{theorem}[Linear overhead at unit conditioning]
    \label{thm:mpu_tree_depth_endpoint}
    \lean{MPUCircuit.linear_overhead_recurrence}
    \leanok
    For every real $C$ and every nonnegative integer $L$,
    \begin{align}
        C(L+1)2^{L+1} & =2(CL2^L)+C2^{L+1}.
    \end{align}
\end{theorem}

\begin{proof}\leanok
    Expand $2^{L+1}=2\cdot2^L$ and distribute multiplication.
\end{proof}

For $N=2^L$, the sequence $T_L=CL2^L$ has order $N\log N$ when $C>0$.
Thus a recurrence with conditioning equal to one need not yield a linear
estimate. This numerical observation is not a lower bound for implementing
an MPU.

\subsection{Boundary Gram matrices of a parity rotation}

The factorization $I_A\otimes cI_B+P_A\otimes isP_B$, with Hermitian
involutions $P_A,P_B$, gives the following two Gram matrices. Their
algebraic inverse-trace identity is distinct from the identification of all
admissible boundary states and the infimum defining the conditioning cost.

\begin{definition}[Parity boundary Gram matrices]
    \label{def:mpu_parity_boundary_grams}
    \lean{MPUCircuit.parityLeftGram, MPUCircuit.parityRightGram}
    \leanok
    For complex numbers $c,s,x,y$, define
    \begin{align}
        G_A(x)     & =\begin{pmatrix}1&x\\x&1\end{pmatrix},            &
        G_B(c,s,y) & =\begin{pmatrix}c^2&icsy\\-icsy&s^2\end{pmatrix}.
    \end{align}
    When the parameters are real, $x,y$ represent the expectations of
    the two involutions in exterior density operators.
\end{definition}

\begin{theorem}[Inverse trace of the parity Gram matrices]
    \label{thm:mpu_parity_inverse_trace}
    \lean{MPUCircuit.trace_inverse_parityGrams}
    \leanok
    For the matrices of Definition~\ref{def:mpu_parity_boundary_grams},
    \begin{align}
        \operatorname{tr}\bigl[G_B(c,s,y)^{-1}(G_A(x)^{-1})^T\bigr]
         & =\frac{c^2+s^2}{c^2s^2(1-x^2)(1-y^2)}.
    \end{align}
    The identity holds for arbitrary complex parameters, with the convention
    that the inverse of a singular matrix and the reciprocal of zero are zero.
\end{theorem}

\begin{proof}\leanok
    \uses{def:mpu_parity_boundary_grams}
    Express each inverse as its adjugate divided by its determinant.
    The determinants are $1-x^2$ and $c^2s^2(1-y^2)$. The mixed
    off-diagonal contributions to the trace cancel, leaving numerator
    $c^2+s^2$.
\end{proof}

\subsection{Interval isometries for product-projection phase gates}

Let $P_j$ be orthogonal projections on the individual sites and $|z|=1$.
The unitary $I+(z-1)\bigotimes_jP_j$ is a product-projection instance of
an MPO-projector unitary from \cite{Styliaris2025MPUCircuits}. The following
boundary frames supply interval isometries directly, without requiring
realization by exterior density operators.

\begin{definition}[Projection interval maps]
    \label{def:mpu_projection_interval_maps}
    \lean{MPUCircuit.projectionPrefixMap, MPUCircuit.projectionInteriorMap, MPUCircuit.projectionSuffixMap}
    \leanok
    For a square complex matrix $P$ and a complex number $z$, define
    \begin{align}
        V_{\mathrm{pre}}(P)   & =\begin{pmatrix}I-P\\P\end{pmatrix},                \\
        V_{\mathrm{int}}(P)   & =\begin{pmatrix}I-P/2\\P/2\\-P/2\\P/2\end{pmatrix}, \\
        V_{\mathrm{suf}}(P,z) & =\begin{pmatrix}I+(z-1)P/2\\(z-1)P/2\end{pmatrix}.
    \end{align}
    The displayed components are stacked vertically. For the interior map,
    their order is the lexicographic order of the two bond labels.
\end{definition}

\begin{theorem}[Normalization of projection interval maps]
    \label{thm:mpu_projection_interval_isometries}
    \lean{MPUCircuit.projectionPrefixMap_isIsometry, MPUCircuit.projectionInteriorMap_isIsometry, MPUCircuit.projectionSuffixMap_isIsometry}
    \leanok
    If $P$ is an orthogonal projection, the prefix and interior maps of
    Definition~\ref{def:mpu_projection_interval_maps} are isometries.
    If also $|z|=1$, the suffix map is an isometry.
\end{theorem}

\begin{proof}\leanok
    \uses{def:mpu_projection_interval_maps}
    Sum the Gram products of the displayed components, using
    $P^\dagger=P$ and $P^2=P$. For the prefix and interior maps this gives
    $I$. For the suffix map, the coefficient of $P$ is
    \begin{align}
        \frac{z-1}{2}+\frac{\bar z-1}{2}+\frac{|z-1|^2}{2}
         & =\frac{|z|^2-1}{2}=0.
    \end{align}
\end{proof}

\begin{definition}[Projection boundary frames and merging kernel]
    \label{def:mpu_projection_boundary_frames}
    \lean{MPUCircuit.projectionLeftFrame, MPUCircuit.projectionRightFrame, MPUCircuit.projectionMergeKernel}
    \leanok
    Define
    \begin{align}
        C_L & =\begin{pmatrix}1&1/2\\0&1/2\end{pmatrix}, &
        C_R & =\begin{pmatrix}1&-1\\0&1\end{pmatrix},    &
        K   & =\begin{pmatrix}1&-1\\1&1\end{pmatrix}.
    \end{align}
    The columns of $C_L,C_R$ are the left and right bond vectors, respectively.
    The matrix $K$ gives the coefficients of their bilinear contraction.
\end{definition}

\begin{theorem}[Projection bond contraction]
    \label{thm:mpu_projection_bond_contraction}
    \lean{MPUCircuit.projectionMergeKernel_contract_frames, MPUCircuit.projectionMergeKernel_trace_conjTranspose_mul, MPUCircuit.projectionFrames_trace_inverse_grams, MPUCircuit.projectionFrames_isUnit}
    \leanok
    The frames in Definition~\ref{def:mpu_projection_boundary_frames} are
    invertible and obey
    \begin{align}
        C_R^T K C_L                      & =I, \\
        \operatorname{tr}(K^\dagger K)   & =4, \\
        \operatorname{tr}\bigl[(C_R^\dagger C_R)^{-1}
        ((C_L^\dagger C_L)^{-1})^T\bigr] & =4.
    \end{align}
\end{theorem}

\begin{proof}\leanok
    \uses{def:mpu_projection_boundary_frames}
    Multiply the two-by-two matrices. Their determinants are
    $\det C_L=1/2$ and $\det C_R=1$. Evaluate the inverse Gram matrices
    by their adjugates and determinants.
\end{proof}

The bond contraction in Theorem~\ref{thm:mpu_projection_bond_contraction}
identifies matching tensor labels. For adjacent intervals it multiplies
$P_{jk}$ and $P_{k+1,l}$ by their tensor product, giving $P_{jl}$.
Its norm is two and is independent of the phase $z$.

\subsection{Finite coefficient alphabets at a cut}

A coefficient matrix at a cut factors through the corresponding MPO bond.
When its coefficients belong to a finite alphabet, the number of distinct
residual coefficient functions can be bounded without a smallest-Schmidt-value
assumption.

\begin{theorem}[Columns distinguishing the rows]
    \label{thm:mpu_cut_row_separator}
    \lean{Matrix.exists_row_separating_columns}
    \leanok
    Let $M$ be a matrix over a field $K$, indexed by an arbitrary row set
    $I$ and a finite column set $J$. There is a subset $J_0\subseteq J$ of
    cardinality $\operatorname{rank}M$ such that two rows of $M$ are equal
    whenever they agree at every column in $J_0$.
\end{theorem}

\begin{proof}\leanok
    Choose a basis of the column space among the columns of $M$.
    Agreement at these columns is preserved under linear combinations,
    so it implies agreement at every column.
\end{proof}

\begin{theorem}[Number of finite-alphabet rows]
    \label{thm:mpu_cut_finite_alphabet_rows}
    \lean{Matrix.card_range_le_pow_rank, Matrix.card_range_le_pow_of_rank_le}
    \leanok
    Let $M$ have an arbitrary row set and a finite column set over a field.
    Suppose every entry belongs to a finite set $S$. The set of distinct
    rows is finite and has cardinality at most $|S|^{\operatorname{rank}M}$.
    If also $S$ is nonempty and $\operatorname{rank}M\leq D$, this number
    is at most $|S|^D$.
\end{theorem}

\begin{proof}\leanok
    \uses{thm:mpu_cut_row_separator}
    Restrict each row to the chosen columns. This injects the set of distinct
    rows into a set of $|S|^{\operatorname{rank}M}$ functions.
    The rank-bound consequence uses $|S|\geq1$.
\end{proof}

\begin{theorem}[Finite-alphabet rows from a bond factorization]
    \label{thm:mpu_cut_bond_row_count}
    \lean{Matrix.card_range_mul_le_pow, Matrix.card_range_mul_le_two_pow}
    \leanok
    Let $A$ and $B$ be matrices over a field with compatible middle index set
    of cardinality $D$ and a finite column set for $B$. Suppose every entry
    of $AB$ belongs to a nonempty finite set $S$. The set of distinct rows
    of $AB$ is finite and has cardinality at most $|S|^D$.
    In particular, zero-one entries give at most $2^D$ distinct rows.
\end{theorem}

\begin{proof}\leanok
    \uses{thm:mpu_cut_finite_alphabet_rows}
    Use $\operatorname{rank}(AB)\leq\operatorname{rank}A\leq D$.
    For the zero-one specialization, take $S=\{0,1\}$.
\end{proof}

\subsection{Unit-circle geometry and diagonal circuits}

The following elementary geometry is the rank-two step in the diagonal
circuit argument of \cite{DiagonalMPUCircuits2026}.

\begin{theorem}[Three points on a vertical line]
    \label{thm:mpu_circle_vertical}
    \lean{Complex.eq_or_eq_or_eq_of_re_eq_of_normSq_eq}
    \leanok
    Let $u,v,w\in\mathbb C$ have the same real part and the same squared
    modulus. Then $u=v$, $u=w$, or $v=w$.
\end{theorem}
\begin{proof}\leanok
    Equality of squared moduli gives equal squares of the imaginary parts.
    Three real numbers with the same square have at most two values.
\end{proof}

\begin{theorem}[Constant real functional on the unit circle]
    \label{thm:mpu_circle_functional}
    \lean{Complex.eq_zero_of_re_mul_eq_of_three_unit_points}
    \leanok
    Let $u,v,w$ be three distinct complex numbers of modulus one. If
    $\operatorname{Re}(zu)=\operatorname{Re}(zv)=\operatorname{Re}(zw)$,
    then $z=0$.
\end{theorem}
\begin{proof}\leanok
    \uses{thm:mpu_circle_vertical}
    The three products have the same real part and the same modulus.
    Two are equal; their factors are distinct, so $z=0$.
\end{proof}

\begin{theorem}[Three affine unit-circle values]
    \label{thm:mpu_circle_affine}
    \lean{Complex.eq_zero_or_eq_zero_of_three_unit_affine_values}
    \leanok
    Let $u,v,w$ be distinct complex numbers of modulus one. If
    $|a+bu|=|a+bv|=|a+bw|=1$, then $a=0$ or $b=0$.
\end{theorem}
\begin{proof}\leanok
    \uses{thm:mpu_circle_functional}
    Expand the squared moduli. The real parts of
    $a\overline b\overline u$, $a\overline b\overline v$ and
    $a\overline b\overline w$ agree. Thus $a\overline b=0$.
\end{proof}

\begin{theorem}[A unit-modulus factorization dichotomy]
    \label{thm:mpu_circle_factorization}
    \lean{Complex.unitCircle_factorization_dichotomy}
    \leanok
    Let $I,J$ be sets with $J$ nonempty. Suppose $a,b:I\to\mathbb C$
    and $t:J\to\mathbb C$ satisfy $|t_j|=1$ and $|a_i+b_it_j|=1$
    for all $i,j$. Then either $a_i=0$ or $b_i=0$ for every $i$, or
    there are $j_0,j_1\in J$ such that $t_j\in\{t_{j_0},t_{j_1}\}$
    for every $j\in J$.
\end{theorem}
\begin{proof}\leanok
    \uses{thm:mpu_circle_affine}
    If $t$ has three distinct values, apply the preceding theorem to each
    row. Otherwise choose one value, and, if it exists, a second distinct
    value. There can be no third.
\end{proof}

% The complete diagonal circuit constructions are not yet formalized.
The construction of \cite{DiagonalMPUCircuits2026} gives an exact circuit with $O(N^{\log_2 3})$ gates for
any diagonal MPU with consecutive operator Schmidt ranks at most two.
A sequential traversal uses $O(N)$ auxiliary qubits; a parallel version
uses $O(N^{\log_2 3})$ auxiliary qubits and logarithmic depth.
For diagonal MPUs of arbitrary maximal bond dimension $D$, reversible
phase evaluation gives operator-norm error at most $\varepsilon$ with
resources polynomial in $N,D,\log(1/\varepsilon)$. These estimates concern diagonal operators.

\begin{theorem}[An affine product of bond factorizations]
    \label{thm:mpu_phase_selector_factorization}
    \lean{Matrix.exists_hadamard_affine_factorization}
    \leanok
    Let $K$ be a field, $P,Q$ finite sets, and let
    $A:I\times P\to K$, $B:P\times J\to K$,
    $C:I\times Q\to K$, and $E:Q\times J\to K$ be matrices.
    For $s,t\in K$, the matrix
    $H_{ij}=s(AB)_{ij}(CE)_{ij}+t$ factors through the index set
    $(P\times Q)\sqcup\{*\}$.
\end{theorem}
\begin{proof}\leanok
    At $(p,q)$ use left entry $sA_{ip}C_{iq}$ and right entry
    $B_{pj}E_{qj}$. At $*$ use left entry $t$ and right entry $1$.
    Distributivity gives their product equal to $H$.
\end{proof}

\begin{theorem}[Boolean selectors from two bond factorizations]
    \label{thm:mpu_phase_selector_row_count}
    \lean{Matrix.card_range_hadamard_affine_le_two_pow}
    \leanok
    Under the preceding hypotheses suppose $J$ is finite and every
    entry of $H$ is zero or one. The set of distinct rows of $H$ is
    finite and has cardinality at most $2^{|P||Q|+1}$.
    In particular, $|P|=|Q|=2$ gives at most 32 rows.
\end{theorem}
\begin{proof}\leanok
    \uses{thm:mpu_phase_selector_factorization, thm:mpu_cut_bond_row_count}
    Apply the finite-alphabet row bound to the displayed factorization.
    Its middle index has cardinality $|P||Q|+1$.
\end{proof}
